% Nội dung báo cáo được chuyển từ bản Markdown đã đối chiếu.
\chapter{NỘI DUNG NGHIÊN CỨU QUA NĂM GIAI ĐOẠN LỊCH SỬ}

\section{Giai đoạn 1954 đến 1960}

\subsection{Bối cảnh lịch sử}

Sau Hiệp định Giơ-ne-vơ năm 1954, miền Bắc được giải phóng, trong khi đất nước tạm thời bị chia cắt. Nhiệm vụ hòa bình, thống nhất chưa hoàn thành. Miền Bắc bước vào khôi phục kinh tế, hàn gắn hậu quả chiến tranh và từng bước xây dựng cơ sở cho chủ nghĩa xã hội. Ở miền Nam, Mỹ từng bước thay thế vai trò của Pháp, hậu thuẫn chính quyền Ngô Đình Diệm, chống lại việc thực hiện tổng tuyển cử thống nhất đất nước. [1, tr. 180 - 183].

Các hoạt động “tố cộng, diệt cộng”, bắt bớ và đàn áp làm lực lượng cách mạng miền Nam chịu nhiều tổn thất. Thực tiễn đặt ra yêu cầu vừa giữ gìn cơ sở, vừa tìm phương thức đấu tranh phù hợp để bảo vệ nhân dân và duy trì phong trào. Chính trong bối cảnh đó, việc điều chỉnh phương pháp đấu tranh trở thành yêu cầu cấp thiết của cách mạng miền Nam. [1, tr. 186 - 192].

\subsection{Chủ trương và đường lối của Đảng}

Sau năm 1954, Đảng xác định nhiệm vụ củng cố miền Bắc, đấu tranh thực hiện Hiệp định Giơ-ne-vơ và giữ gìn lực lượng cách mạng ở miền Nam. Trong điều kiện mới, yêu cầu bảo vệ hòa bình gắn với việc chuẩn bị lực lượng lâu dài, không tách nhiệm vụ của từng miền khỏi mục tiêu chung của cả nước.

Hội nghị Trung ương 15 năm 1959 là bước chuyển quan trọng về phương pháp cách mạng miền Nam. Trung ương xác định con đường cơ bản là khởi nghĩa giành chính quyền về tay nhân dân, dựa vào sức mạnh chính trị của quần chúng, kết hợp với lực lượng vũ trang. Chủ trương này đáp ứng yêu cầu bức thiết của phong trào, mở ra khả năng kết hợp đấu tranh chính trị với đấu tranh vũ trang thay vì chỉ dựa vào một hình thức đơn lẻ. [1, tr. 190 - 191].

Đại hội đại biểu toàn quốc lần thứ III của Đảng tháng 9/1960 xác định tiến hành đồng thời cách mạng xã hội chủ nghĩa ở miền Bắc và cách mạng dân tộc dân chủ nhân dân ở miền Nam. Miền Bắc giữ vai trò quyết định nhất đối với sự phát triển của cách mạng cả nước; miền Nam giữ vai trò quyết định trực tiếp đối với sự nghiệp giải phóng miền Nam. Hai nhiệm vụ có quan hệ chặt chẽ, cùng hướng tới hòa bình và thống nhất Tổ quốc. Đây là cơ sở chiến lược để xây dựng hậu phương, phát triển tiền tuyến và huy động sức mạnh cả nước trong những giai đoạn tiếp theo. [1, tr. 193 - 195].

\subsection{Kết quả thực tiễn và ý nghĩa}

Phong trào Đồng Khởi năm 1959 - 1960, tiêu biểu tại Bến Tre ngày 17/1/1960, thể hiện sự phát triển của phong trào quần chúng khi đấu tranh chính trị kết hợp với đấu tranh vũ trang. Những cuộc nổi dậy làm suy yếu bộ máy kiểm soát ở nhiều vùng nông thôn, mở rộng quyền làm chủ của nhân dân và tạo bước chuyển từ thế giữ gìn lực lượng sang thế tiến công. Vai trò của phụ nữ và lực lượng đấu tranh chính trị cho thấy sức mạnh cách mạng không chỉ nằm ở hoạt động quân sự.

Ngày 20/12/1960, Mặt trận Dân tộc Giải phóng miền Nam Việt Nam ra đời, tạo hình thức tổ chức để tập hợp rộng rãi các lực lượng yêu nước. Những chuyển biến này cho thấy mối liên hệ giữa chủ trương đúng với khả năng tổ chức nhân dân thực hiện; đồng thời tạo tiền đề để cách mạng miền Nam đối phó với chiến lược “Chiến tranh đặc biệt” từ năm 1961. [1, tr. 191 - 192].

\subsection{Liên hệ chứng tích lịch sử}

Chứng tích máy chém phản ánh một phương diện của chính sách đàn áp và hoàn cảnh hiểm nghèo mà cán bộ, nhân dân miền Nam phải đối diện. Khi đặt bên cạnh tư liệu về Đồng Khởi và Đội quân tóc dài, người học thấy rõ hơn vai trò của lực lượng chính trị quần chúng, đặc biệt là sự tham gia của phụ nữ trong đấu tranh giành quyền làm chủ. Mối liên hệ giữa bối cảnh đàn áp, yêu cầu chuyển hướng đấu tranh và sự phát triển của phong trào quần chúng góp phần làm sáng tỏ ý nghĩa của Nghị quyết Trung ương 15.

\section{Giai đoạn 1961 đến 1965}

\subsection{Bối cảnh lịch sử}

Từ năm 1961, Mỹ triển khai chiến lược “Chiến tranh đặc biệt”, sử dụng quân đội Sài Gòn làm lực lượng chủ yếu, dưới sự hỗ trợ của cố vấn, vũ khí và phương tiện Mỹ. Hệ thống ấp chiến lược được xây dựng nhằm kiểm soát dân cư, tách quần chúng khỏi lực lượng cách mạng. Những chiến thuật cơ động như trực thăng vận và thiết xa vận tạo thêm khó khăn cho lực lượng vũ trang và phong trào tại chỗ.

Vấn đề đặt ra không chỉ là đối phó với ưu thế trang bị, mà còn là giữ được cơ sở trong nhân dân, phá thế kiểm soát ở nông thôn và duy trì hoạt động trên nhiều địa bàn. Nếu tách đấu tranh quân sự khỏi đấu tranh chính trị và công tác vận động quần chúng, cách mạng khó tạo được sức mạnh đủ rộng để đối phó với chiến lược mới. [1, tr. 200 - 201].

\subsection{Chủ trương và đường lối của Đảng}

Trên cơ sở thế tiến công sau Đồng Khởi, Đảng chủ trương phát triển đấu tranh vũ trang song song với đấu tranh chính trị, đưa cách mạng miền Nam từ khởi nghĩa từng phần lên chiến tranh cách mạng. Việc thống nhất lực lượng vũ trang thành Quân Giải phóng miền Nam Việt Nam năm 1961 và tăng cường tổ chức lãnh đạo đáp ứng yêu cầu phát triển của cuộc đấu tranh.

Phương châm kết hợp quân sự, chính trị và binh vận được triển khai trên ba vùng chiến lược: rừng núi, nông thôn đồng bằng và đô thị. Mỗi địa bàn có điều kiện khác nhau nên cần lựa chọn cách đấu tranh thích hợp, nhưng vẫn hỗ trợ mục tiêu chung. Ở nông thôn, hoạt động quân sự phải gắn với chống dồn dân, phá ấp và giành quyền làm chủ; ở đô thị, phong trào chính trị góp phần gây sức ép và làm suy yếu cơ sở xã hội của chính quyền Sài Gòn.

Sự phối hợp đó thể hiện tư duy chiến tranh nhân dân: không đối đầu với ưu thế vật chất của đối phương chỉ bằng một cách đánh, mà tạo thế đấu tranh rộng khắp, kết hợp lực lượng và hình thức phù hợp với thực tiễn. Hậu phương miền Bắc tiếp tục chi viện để củng cố khả năng chiến đấu của miền Nam. [1, tr. 201 - 202].

\subsection{Kết quả thực tiễn và ý nghĩa}

Chiến thắng Ấp Bắc ngày 2/1/1963 cho thấy lực lượng cách mạng có thể chống lại những cuộc hành quân sử dụng phương tiện cơ động và hỏa lực mạnh. Thắng lợi góp phần củng cố niềm tin, phát triển kinh nghiệm tác chiến và thúc đẩy phong trào thi đua trên chiến trường. Cùng với đấu tranh quân sự, phong trào phá ấp chiến lược và đấu tranh chính trị tại đô thị tạo thêm sức ép đối với chính quyền Sài Gòn.

Các thắng lợi Bình Giã, Ba Gia, Đồng Xoài trong những năm 1964 - 1965 đánh dấu bước phát triển của lực lượng vũ trang và khả năng tổ chức những trận đánh quy mô lớn hơn. Đến giữa năm 1965, chiến lược “Chiến tranh đặc biệt” bị đánh bại. Kết quả này không thể quy về một trận đánh riêng lẻ, mà là sản phẩm của quá trình kết hợp các mặt đấu tranh, bảo vệ cơ sở quần chúng và phát huy sự chi viện từ miền Bắc. [1, tr. 202 - 205].

\subsection{Liên hệ chứng tích lịch sử}

Ảnh tư liệu về ấp chiến lược và hiện vật trực thăng UH-1 gợi lên hai phương diện của chiến lược “Chiến tranh đặc biệt”: kiểm soát dân cư và cơ động quân sự. Hệ thống rào chắn, dồn dân giúp lý giải vì sao giành dân, giữ đất là nhiệm vụ quan trọng; phương tiện không vận giúp hình dung những khó khăn về tương quan trang bị trên chiến trường. Qua đó, có thể thấy ý nghĩa của việc kết hợp đấu tranh quân sự với đấu tranh chính trị và binh vận, vừa đối phó các cuộc hành quân, vừa bảo vệ cơ sở cách mạng trong nhân dân.

\section{Giai đoạn 1965 đến 1968}

\subsection{Bối cảnh lịch sử}

Sau thất bại của “Chiến tranh đặc biệt”, Mỹ chuyển sang “Chiến tranh cục bộ”, trực tiếp đưa quân viễn chinh và quân đồng minh vào tham chiến quy mô lớn ở miền Nam. Đồng thời, Mỹ mở rộng chiến tranh phá hoại miền Bắc bằng không quân và hải quân nhằm làm suy giảm tiềm lực hậu phương và ngăn chặn chi viện.

Chiến tranh mở rộng trên cả hai miền, đặt cách mạng Việt Nam trước thử thách lớn về quân sự, kinh tế và đời sống xã hội. Miền Nam phải đối phó với các cuộc hành quân tìm diệt, bình định; miền Bắc vừa bảo vệ cơ sở sản xuất, giao thông, dân cư, vừa duy trì vai trò hậu phương. Yêu cầu đặt ra là giữ vững quyết tâm, đánh giá đúng khả năng của đối phương và xây dựng cách đánh phù hợp với điều kiện chiến tranh mới. [1, tr. 206 - 208].

\subsection{Chủ trương và đường lối của Đảng}

Hội nghị Trung ương 11 và Hội nghị Trung ương 12 năm 1965 phát triển đường lối kháng chiến trong hoàn cảnh cả nước có chiến tranh. Đảng khẳng định quyết tâm chống Mỹ, cứu nước, phát huy sức mạnh chiến tranh nhân dân và kết hợp chặt chẽ nhiệm vụ ở hai miền. Miền Bắc không chỉ là nơi cung cấp nguồn lực mà còn trực tiếp chiến đấu bảo vệ hậu phương; miền Nam tiếp tục giữ vai trò tiền tuyến lớn.

Ở miền Nam, chủ trương tiến công gắn với phối hợp đấu tranh quân sự, chính trị và hoạt động của quần chúng trên các địa bàn. Mục tiêu không phải chỉ tiêu hao lực lượng quân sự đối phương, mà còn làm thất bại mục đích chiến tranh và ý chí tiếp tục xâm lược. Lời kêu gọi của Chủ tịch Hồ Chí Minh ngày 17/7/1966 tiếp tục khẳng định quyết tâm bảo vệ độc lập, tự do của dân tộc.

Trên cơ sở đánh giá tình hình sau hai mùa khô, Bộ Chính trị và Hội nghị Trung ương 14 tháng 1/1968 quyết định mở cuộc Tổng tiến công và nổi dậy Tết Mậu Thân. Việc lựa chọn đô thị và các cơ quan trọng yếu làm hướng tiến công chủ yếu nhằm tạo chuyển biến lớn, tác động đồng thời tới cục diện quân sự, chính trị và ngoại giao. [1, tr. 207 - 208], [1, tr. 216 - 217].

\subsection{Kết quả thực tiễn và những vấn đề đặt ra}

Các trận Núi Thành, Vạn Tường năm 1965 và kết quả chiến đấu qua hai mùa khô 1965 - 1966, 1966 - 1967 thể hiện khả năng đánh quân Mỹ của quân và dân miền Nam. Những kinh nghiệm tích lũy trong thực tiễn giúp hạn chế một phần ưu thế hỏa lực, cơ động của đối phương và duy trì thế đấu tranh lâu dài. Miền Bắc đồng thời chiến đấu chống chiến tranh phá hoại, bảo vệ các tuyến chi viện và tiếp tục cung cấp nguồn lực cho tiền tuyến.

Cuộc Tổng tiến công và nổi dậy Tết Mậu Thân 1968 gây chấn động lớn, tác động tới dư luận và đời sống chính trị nước Mỹ. Sự kiện góp phần làm thay đổi cách đánh giá về khả năng kết thúc chiến tranh bằng sức mạnh quân sự, thúc đẩy quá trình Mỹ xuống thang và tham gia đàm phán tại Paris. Đây là bước ngoặt quan trọng trong việc kết hợp đấu tranh trên chiến trường với mặt trận ngoại giao.

Đánh giá Mậu Thân cần nhìn nhận cả kết quả chiến lược lẫn tổn thất. Mục tiêu giành thắng lợi quyết định thông qua tổng khởi nghĩa ở đô thị không đạt như dự kiến; lực lượng và cơ sở cách mạng chịu tổn thất nặng. Nhận thức đầy đủ cả hai mặt giúp tránh đồng nhất ý nghĩa của bước ngoặt với việc mọi mục tiêu ban đầu đều thành công, đồng thời rút ra bài học về đánh giá tương quan lực lượng, xác định mục tiêu và tổ chức thực hiện. [1, tr. 214 - 219].

\subsection{Liên hệ chứng tích lịch sử}

Tư liệu ảnh về nhà ga Nam Định bị ném bom cho thấy chiến tranh phá hoại tác động trực tiếp đến hạ tầng giao thông và đời sống dân cư miền Bắc. Việc bảo vệ, khôi phục các tuyến vận tải vì vậy gắn liền với cả nhiệm vụ duy trì sản xuất và bảo đảm chi viện cho tiền tuyến. Những hình ảnh này làm rõ hơn mối quan hệ giữa xây dựng, bảo vệ hậu phương và phát triển cuộc kháng chiến ở miền Nam.

Hiện vật máy bay A-37 bổ sung góc nhìn về phương tiện và hỏa lực không quân trong chiến tranh. Đặt bên cạnh các tư liệu về hạ tầng bị tàn phá, hiện vật giúp hình dung quy mô phương tiện chiến tranh mà nhân dân Việt Nam phải đối mặt, đồng thời gợi nhận thức về sự hy sinh trong quá trình bảo vệ hậu phương và giữ vững sức chi viện. [2].

\section{Giai đoạn 1969 đến 1973}

\subsection{Bối cảnh lịch sử}

Từ năm 1969, chính quyền Tổng thống Richard Nixon triển khai chiến lược “Việt Nam hóa chiến tranh”, từng bước rút quân chiến đấu Mỹ nhưng tiếp tục tăng cường cho quân đội Sài Gòn và duy trì sự hỗ trợ về hỏa lực, trang bị. Việc rút quân Mỹ vì vậy không đồng nghĩa chiến tranh tự chấm dứt. Đối phương tiếp tục bình định, mở rộng hoạt động quân sự và tìm cách ngăn chặn nguồn chi viện cho cách mạng miền Nam.

Chiến tranh còn mở rộng sang Lào và Campuchia, trong khi quan hệ giữa các nước lớn có những thay đổi phức tạp. Sau khi Chủ tịch Hồ Chí Minh qua đời năm 1969, toàn Đảng, toàn dân tiếp tục nhiệm vụ kháng chiến và thực hiện Di chúc của Người. Điều kiện mới đòi hỏi vừa phát huy nội lực, vừa giữ vững độc lập, tự chủ trong xử lý quan hệ quốc tế, đồng thời tăng cường phối hợp giữa các mặt trận. [1, tr. 220 - 221], [1, tr. 223 - 226].

\subsection{Chủ trương và đường lối của Đảng}

Đảng chủ trương giữ vững chiến lược tiến công, chống bình định, củng cố lực lượng và khôi phục, phát triển khả năng chi viện của miền Bắc. Đấu tranh quân sự, chính trị và ngoại giao được kết hợp theo phương châm vừa đánh, vừa đàm. Trong đó, kết quả trên chiến trường tạo điều kiện cho đàm phán; đấu tranh ngoại giao góp phần chuyển hóa kết quả quân sự, chính trị thành những cam kết và điều kiện có lợi cho cách mạng.

Trước việc chiến tranh mở rộng ở Đông Dương, sự phối hợp với nhân dân Lào và Campuchia có ý nghĩa quan trọng trong bảo vệ các tuyến chiến lược và chống lại những cuộc hành quân của đối phương. Sự đoàn kết này không thay thế vai trò tự lực của nhân dân Việt Nam, mà bổ sung sức mạnh trong một không gian chiến tranh có liên hệ chặt chẽ.

Ở miền Bắc, nhiệm vụ khôi phục kinh tế, bảo đảm giao thông vận tải và chuẩn bị chiến đấu tiếp tục gắn với trách nhiệm hậu phương. Chủ trương đó cho thấy đấu tranh giành thắng lợi ở miền Nam không tách rời năng lực bảo vệ và duy trì nguồn lực của cả nước. [1, tr. 220 - 226].

\subsection{Kết quả thực tiễn và ý nghĩa}

Chiến thắng Đường 9 - Nam Lào năm 1971 cùng những kết quả phối hợp chiến đấu ở Đông Dương làm thất bại các nỗ lực cắt đứt tuyến chi viện và bộc lộ hạn chế của chiến lược “Việt Nam hóa chiến tranh”. Cuộc tiến công chiến lược năm 1972 tiếp tục tác động mạnh tới cục diện, buộc Mỹ tăng cường sử dụng không quân, hải quân và mở lại chiến tranh phá hoại miền Bắc.

Trong tháng 12/1972, quân dân miền Bắc đánh trả cuộc tập kích đường không chiến lược của Mỹ vào Hà Nội, Hải Phòng. Cùng với quá trình đấu tranh quân sự, chính trị và ngoại giao trước đó, kết quả này góp phần tạo điều kiện đi tới ký kết Hiệp định Paris ngày 27/1/1973. Hiệp định ghi nhận các quyền dân tộc cơ bản của Việt Nam và quy định việc rút quân Mỹ cùng quân đồng minh, tạo thay đổi quan trọng về tương quan lực lượng.

Hiệp định Paris là thắng lợi có ý nghĩa chiến lược, nhưng không phải thời điểm toàn bộ chiến tranh kết thúc. Xung đột tiếp tục diễn ra ở miền Nam sau năm 1973, đặt ra nhiệm vụ bảo vệ thành quả, củng cố lực lượng và hoàn thành sự nghiệp giải phóng. Nhìn đúng giới hạn và ý nghĩa của Hiệp định giúp thấy rõ quan hệ giữa hai bước phát triển: buộc quân Mỹ rút khỏi Việt Nam và tiến tới thắng lợi hoàn toàn năm 1975. [1, tr. 227 - 229].

\subsection{Liên hệ chứng tích lịch sử}

Tư liệu về hậu quả chất độc da cam và hình ảnh bà Nguyễn Thị Bình ký Hiệp định Paris thể hiện hai phương diện của cuộc chiến: những tổn thương đối với con người và nỗ lực đấu tranh ngoại giao nhằm lập lại hòa bình. Hậu quả chất độc da cam kéo dài qua nhiều thời kỳ, không chỉ giới hạn ở các năm 1969 - 1973, cho thấy chiến tranh tiếp tục tác động đến đời sống ngay cả khi tiếng súng đã chấm dứt.

Hình ảnh lễ ký Hiệp định góp phần làm rõ vai trò của mặt trận ngoại giao và sự tham gia của phụ nữ Việt Nam trong cuộc đấu tranh vì độc lập, hòa bình. Những tư liệu về nạn nhân chiến tranh đồng thời gợi lên sự cảm thông và trách nhiệm hỗ trợ những người chịu hậu quả lâu dài. Từ đó, việc tìm hiểu lịch sử không chỉ bồi dưỡng nhận thức về thắng lợi mà còn củng cố ý thức nhân ái và trách nhiệm gìn giữ hòa bình.

\section{Giai đoạn 1973 đến 1975}

\subsection{Bối cảnh lịch sử}

Sau Hiệp định Paris, quân Mỹ và quân đồng minh rút khỏi Việt Nam, tạo điều kiện thuận lợi hơn cho cách mạng miền Nam. Tuy nhiên, chính quyền Sài Gòn tiếp tục các hoạt động quân sự và lấn chiếm. Thành quả đạt được chưa đồng nghĩa hòa bình đã được bảo đảm; yêu cầu bảo vệ lực lượng, vùng giải phóng và quyền làm chủ của nhân dân vẫn rất cấp thiết.

Trong những năm 1973 - 1974, miền Bắc khôi phục kinh tế, tăng cường chi viện; hệ thống giao thông chiến lược và hậu cần tiếp tục phát triển. Việc xây dựng các quân đoàn chủ lực cơ động tạo thêm khả năng tác chiến hiệp đồng quy mô lớn. Những điều kiện đó tích lũy thế và lực cho giai đoạn quyết định, thay vì coi thắng lợi năm 1975 là kết quả của một thời cơ xuất hiện tách rời quá trình chuẩn bị. [1, tr. 228 - 231].

\subsection{Chủ trương và đường lối của Đảng}

Hội nghị Trung ương 21 năm 1973 khẳng định cách mạng miền Nam tiếp tục con đường bạo lực cách mạng, giữ vững chiến lược tiến công và chủ động phát triển thực lực. Chủ trương này định hướng việc chống lấn chiếm, bảo vệ thành quả và chuẩn bị tiến lên giải phóng hoàn toàn miền Nam, không thụ động trông chờ vào việc đối phương thực hiện Hiệp định.

Cuối năm 1974 và đầu năm 1975, Bộ Chính trị họp hai đợt để bàn kế hoạch giải phóng miền Nam trong hai năm 1975 - 1976. Đồng thời, kế hoạch dự kiến khả năng hoàn thành ngay trong năm 1975 nếu thời cơ xuất hiện. Điều này thể hiện sự kết hợp giữa chuẩn bị có hệ thống với khả năng điều chỉnh theo diễn biến thực tế.

Thắng lợi Phước Long ngày 6/1/1975 cung cấp thêm cơ sở đánh giá tương quan lực lượng và phản ứng của Mỹ. Khi cuộc tiến công ở Tây Nguyên phát triển thuận lợi, ngày 18/3/1975, Bộ Chính trị quyết định giải phóng miền Nam trong năm 1975; ngày 25/3 tiếp tục bổ sung quyết tâm hoàn thành trước mùa mưa. Những quyết định nối tiếp cho thấy thời cơ được nhận diện và khai thác trong quá trình chiến đấu, gắn với năng lực tập trung lực lượng và tổ chức thực hiện. [1, tr. 230 - 232].

\subsection{Kết quả thực tiễn và ý nghĩa}

Cuộc Tổng tiến công và nổi dậy mùa Xuân 1975 phát triển qua những đòn tiến công chiến lược lớn. Chiến thắng Buôn Ma Thuột ngày 10/3 mở ra sự chuyển biến nhanh chóng ở Tây Nguyên. Việc giải phóng Huế ngày 26/3 và Đà Nẵng ngày 29/3 tiếp tục làm thay đổi cục diện, tạo điều kiện tập trung lực lượng cho hướng Sài Gòn - Gia Định.

Chiến dịch Hồ Chí Minh bắt đầu ngày 26/4/1975. Trưa ngày 30/4, chính quyền Sài Gòn tuyên bố đầu hàng, Sài Gòn được giải phóng; các địa phương còn lại tiếp tục được giải phóng sau đó. Thắng lợi là kết quả của sự phối hợp giữa các binh đoàn chủ lực, lực lượng tại chỗ, hoạt động nổi dậy của quần chúng và sự bảo đảm hậu cần từ quá trình chuẩn bị lâu dài.

Đại thắng mùa Xuân năm 1975 kết thúc thắng lợi cuộc kháng chiến chống Mỹ, cứu nước, giải phóng hoàn toàn miền Nam và tạo điều kiện quyết định cho việc hoàn thành thống nhất đất nước về mặt nhà nước. Cần phân biệt việc giải phóng hoàn toàn miền Nam, thống nhất đất nước trên thực tế năm 1975 với tiến trình thống nhất về mặt nhà nước được hoàn thành năm 1976. [1, tr. 232 - 233], [1, tr. 238 - 241].

\subsection{Liên hệ chứng tích lịch sử}

Các hiện vật pháo M107, xe tăng M48 và máy bay A-37 phản ánh quy mô trang bị của các lực lượng tham chiến. Khi liên hệ với tư liệu về tuyến vận tải Trường Sơn, có thể thấy vai trò của công tác hậu cần, khả năng cơ động và sự chuẩn bị lực lượng đối với các chiến dịch lớn. Sự phát triển của thế trận và nguồn lực là cơ sở quan trọng để quân, dân ta tận dụng thời cơ trong mùa Xuân năm 1975.

Tư liệu về Dinh Độc Lập và sự kiện ngày 30/4/1975 bổ sung hình dung về thời điểm kết thúc chiến dịch giải phóng Sài Gòn. Đặt trong toàn bộ tiến trình kháng chiến, những hình ảnh ấy không chỉ biểu hiện thắng lợi của một chiến dịch mà còn gợi nhắc quá trình đấu tranh, chuẩn bị và hy sinh lâu dài của nhân dân cả nước. [1, tr. 230 - 233].

Qua năm giai đoạn, có thể thấy sự phát triển liên tục của quá trình lãnh đạo: giữ gìn và xây dựng lực lượng; phát triển chiến tranh cách mạng; đối phó với chiến tranh mở rộng; kết hợp đấu tranh để buộc quân Mỹ rút; và chủ động tạo thế, nắm thời cơ để giành thắng lợi cuối cùng. Mỗi giai đoạn có nhiệm vụ và phương thức riêng, nhưng cùng gắn với mục tiêu độc lập dân tộc và thống nhất đất nước.

\clearpage
